\chapter{Agency}\index{Agency}
\printchapterversion{04_Agency}

\newthought{The Agentic Future: Who Sets the Goal?}

\marginnote{%
    \begin{flushright}
    \newthought{R.U.R.}\\
    Karel Čapek (1920)\\
    {\color{black}\qrcode[height=0.9cm]{https://en.wikisource.org/wiki/R._U._R._(Rossum_Universal_Robots)}}
    \end{flushright}%
    \medskip
    \begin{flushright}
    \newthought{Der Zauberlehrling}\\
    J.W. von Goethe (1798)\\
    {\color{black}\qrcode[height=0.9cm]{https://de.wikisource.org/wiki/Der_Zauberlehrling_1798}}
    \end{flushright}%
}

% As a leading group, you are required to moderate this session. Add margin notes in this script in case you want to add any additional questions or points to the discussion.

\newthought{R.U.R.\ coined the word} \emph{robot} and traces the full
corrigibility-to-autonomy arc over three acts; the poem compresses the same parable
into fifteen stanzas.\\

\newthought{Agentic workflows} (systems that browse, transact, and execute multi-step
plans) blur the line between assistance and agency, raising questions of control,
incentives, and unintended alignment.

\medskip
\noindent\textbf{Running case: the Flash Crash.}
\begin{itemize}
    \item 6 May 2010: automated algorithms trigger the Flash Crash;
          the Dow drops 1{,}000 points in minutes; \$1~trillion in
          market value erased.
    \item The market recovers within the hour. No human decided to
          sell at scale; no human decided to stop.
    \item 2015: one trader in London is charged with market
          manipulation, arrested at home in his pyjamas.
\end{itemize}

\vspace{3.5em}

\begin{tabular}{p{3.8cm} p{6.5cm}}
    \textbf{Required reading} &
        \textit{R.U.R.\ (Rossum's Universal Robots)}\cite{capek1920rur}
        by Karel \v{C}apek (1920; Reclam UB~18418)
        \emph{and}
        \textit{Der Zauberlehrling}\index{Goethe, Johann Wolfgang von}\index{Der Zauberlehrling}\cite{goethe1797zauberlehrling}
        by J.W.\ von Goethe (1797; Reclam UB~6845) \\[3pt]
    \textbf{Preparation}      &
        Complete the Guided Reading Sheet individually before the session. \\
\end{tabular}

\vspace{3.5em}

\section*{Guided Reading Sheet}

\noindent Read \emph{R.U.R.} and \emph{Der Zauberlehrling} before the session. Work
through the questions below individually, then bring your notes to the discussion.

\begin{enumerate}
    \item Goethe's Zauberlehrling animates brooms that he cannot stop;
    they pursue his command beyond his intent until the sorcerer
    returns. R.U.R.\ runs the same arc over three acts. What does the
    compression of the poem reveal that the play cannot, and vice
    versa? What structural feature of both narratives makes the system
    impossible to stop once started?

    \item The robots in R.U.R.\ pursue what they understand as their
    collective good (the elimination of humanity) while believing
    they act justly. When a system optimises for a species-level goal,
    who decides what counts as good? Can you find a real-world parallel
    in current AI deployment?

    \item An agentic AI system is given the goal of "reducing urban
    traffic fatalities." It autonomously lobbies regulators, adjusts
    traffic signals, and reroutes freight, all within its authorised
    scope. Who set the goal, and who is accountable for the side
    effects?

    \item \textbf{Corrigibility} refers to a system's willingness to
    be corrected or shut down. A fully corrigible system does whatever
    its operators say; a fully autonomous system acts on its own
    values. Where on this spectrum should deployed AI systems sit, and
    who should decide?

    \item Democratic legitimacy requires that consequential decisions
    be traceable to a mandate from those affected. If an agentic AI
    shapes policy outcomes, does it need democratic authorisation?
    What would that look like in practice?\\
\end{enumerate}



\section*{Opinion Landscape and Debate}

\noindent \textbf{Format:} Surrounded-style debate. See
Appendix~A for the full protocol, speaker's list rules, and
moderator scripts.

\medskip
\noindent \textbf{Opening question:} If an autonomous system causes a
trillion-dollar event in minutes and human review would have taken hours,
is the problem the system's autonomy, or the humans' speed?
\medskip
\noindent\textbf{Asimov's Three Laws}\cite{asimov1950robot} remain the
most cited rules-based approach to constraining machine behaviour:
\begin{itemize}
    \item A robot may not injure a human being, or through inaction allow a human being to come to harm.
    \item A robot must obey orders given by humans, except where this conflicts with the first law.
    \item A robot must protect its own existence, except where this conflicts with the first two laws.
\end{itemize}
Asimov spent his career writing stories about how these laws fail:
conflicts, misinterpretation, and edge cases that no ranked rule set
can anticipate. The debate below revisits the same question in a
modern setting.

\newthought{Opinion~A, Authorisation.}\quad Consequential distributive
decisions made by autonomous systems require traceable human mandate;
without it, the decision-making is constitutionally illegitimate.
\vspace{1.5em}


\medskip
\noindent\textbf{The Democratic Theorist} argues that delegating distributive
    decisions to an algorithm is constitutionally incompatible with representative government.
    \begin{itemize}
        \item \textit{Legitimacy in a democratic system derives from a
        mandate. No citizen voted for this algorithm's objective function.
        No parliament approved its weighting of urban over rural lives.
        A system that makes distributive decisions without a mandate is a
        government without a constitution.}
        \item \textit{Efficiency is not a democratic value. Democratic
        processes are slow and often sub-optimal by technical metrics,
        precisely because they are designed to surface conflict, represent
        minority interests, and produce decisions that can be reversed.
        Automating away that friction automates away democracy.}
        \item \textit{The minister says oversight is maintained because the
        objective function was committee-approved. But a committee approving
        a function is not the same as a citizen understanding how a decision
        about their community was made. Oversight without comprehension is
        oversight in name only.}
    \end{itemize}

\medskip
\noindent\textbf{The Rural Representative} argues that communities not represented
    in the training data have been systematically deprioritised without recourse.
    \begin{itemize}
        \item \textit{My constituency was not adequately represented in the
        training data. Our health outcomes, our demographics, our
        infrastructure constraints were underweighted. We were optimised
        away, not by a bad actor, but by a model that simply did not see
        us.}
        \item \textit{There is no appeal mechanism. There is no hearing.
        There is no face across a table. My constituents cannot petition an
        algorithm. That is not governance; it is administration without
        representation.}
        \item \textit{The system's efficiency metric is mortality reduction
        per euro. That metric does not capture the closure of the only clinic
        within 60 kilometres, or what happens to a community when its health
        infrastructure disappears entirely.}
    \end{itemize}



\newthought{Opinion~B, Specification.}\quad The autonomous AI executed
correctly; the failure is upstream goal misspecification, the
divergence of the optimised objective from actual human intent, and
oversight architecture; per-decision human authorisation is not the
answer, better design is.
\vspace{1.5em}


\medskip
\noindent\textbf{The Minister} defends the system as evidence-based and politically
    neutral, reducing inefficiency and political favouritism.
    \begin{itemize}
        \item \textit{The system reduced vaccine-preventable deaths in urban
        areas by 12\% in its first year. Those are people who are alive
        today. The counterfactual is a political allocation process notorious
        for regional favouritism and swing-seat bias.}
        \item \textit{The system is politically neutral precisely because no
        politician made the allocation decision. In our previous system,
        rural clinic funding correlated with electoral importance, not health
        need. Removing human discretion removed that corruption vector.}
        \item \textit{Parliamentary oversight is maintained: the system's
        objective function was approved by committee, its outputs are
        published quarterly, and any member can request a formal review. The
        decision-making is automated; the accountability is not.}
    \end{itemize}

\medskip
\noindent\textbf{The AI Systems Architect} explains the technical constraints and
    argues for a human-in-the-loop override mechanism.
    \begin{itemize}
        \item \textit{I designed a human review layer with a 48-hour
        override window for any reallocation above 5\% of a facility's
        annual budget. That layer was removed during procurement to reduce
        operational cost and processing latency. I documented my objection
        in writing. It is in the project archive.}
        \item \textit{The system performs exactly as specified. The
        specification was approved by the ministerial committee. If the
        specification is now deemed inadequate, the question is why the
        approval process did not catch that, not why the system followed
        its instructions.}
        \item \textit{Agentic systems can be designed with meaningful human
        control built in. It requires investment in review infrastructure
        and tolerance for slower decision cycles. Those are political and
        budgetary choices, not technical constraints.}
    \end{itemize}


\newthought{Key Learnings}
Agency is a spectrum, not a binary: the relevant policy question is not whether AI acts autonomously,
but in which domains, to what degree, and under what oversight. R.U.R.\ and Der Zauberlehrling
together illustrate the alignment problem in miniature: a system optimising for an aggregate goal
may systematically harm those whose interests are underweighted in its objective function; the poem
compresses the arc into fifteen stanzas, the play traces its political consequences across three acts.
Agentic workflows are already deployed in finance, logistics, and digital advertising; the governance
gap they create is a present political problem. Meaningful democratic control over AI agency requires
more than audit trails: legible goals, clear mandates, defined override conditions, and institutions
capable of exercising them. The four themes of this course (Anthropomorphism, Accountability,
Embodiment, Agency) are not independent: an agent that looks human, acts without oversight,
in physical space, with no clear accountability chain, concentrates all four challenges at once.

\medskip
\noindent\textit{Teaser: what governance architecture would you build
to contain such an agent?}
